\section{Motivation: the missing dimension}
\label{sec:motivation}

Modern knowledge graph (KG) construction has converged on a recognizable pattern. Documents enter; an extractor (whether rule-based, fine-tuned NER, or LLM-driven) emits entities and typed relations; a graph builder deduplicates and resolves; a visualizer renders. The output is a labeled property graph that answers structural questions: which entities are most central, which clusters form, which relations are most prevalent. The flagship knowledge-graph platforms -- Neo4j Bloom, Stardog, GraphDB, Microsoft's GraphRAG~\cite{edge2024graphrag}, and a growing ecosystem of open-source LLM-based extractors -- all deliver this primary product.

A flat knowledge graph, however, is fundamentally lossy with respect to the epistemic content of its source corpus. Consider three sentences that might appear in a drug-discovery corpus:

\begin{enumerate}
\item ``In SURPASS-2, tirzepatide 15 mg reduced HbA1c by 2.30\% versus 1.86\% for semaglutide 1 mg (p<0.001).'' \cite{frias2021surpass2}
\item ``In one or more embodiments, the compositions of the present invention may be useful for the treatment of cardiovascular disease, neurodegenerative disorders, and metabolic dysfunction.''
\item ``GLP-1 receptor agonism has been hypothesized to confer neuroprotective effects through mechanisms that are not yet fully characterized.''
\end{enumerate}

A flat KG would extract three relations of the form $\langle\text{compound}\rangle$ \texttt{INDICATED\_FOR} $\langle\text{disease}\rangle$ from these sentences and present them as equivalent edges in the graph. Yet they are not equivalent. The first is an \emph{asserted} relation backed by a randomized controlled trial. The second is \emph{prophetic} -- patent forward-looking language designed to maximize claim breadth, with no requirement that the indication be empirically established. The third is \emph{hypothesized} -- explicitly hedged research speculation. A drug-discovery analyst reading the graph cannot make decisions on the basis of edges that flatten this distinction.

This is not a problem that domain-specific extraction prompting solves. The information is present in the source text -- the syntax of patents, hedging vocabulary in research papers, the structural conventions of trial reports -- but a graph that records only $\langle$ subject, predicate, object, confidence, source $\rangle$ has no field in which to store it. The missing dimension is \emph{epistemic typing}: a categorical label, attached to each relation, that records what kind of claim it is.

This paper introduces Epistract, a domain-agnostic knowledge graph framework whose primary contribution is to treat the epistemic dimension as a first-class architectural layer rather than as ad-hoc per-relation metadata. Epistract produces \emph{two-layer} graphs: Layer 1 is the conventional brute-facts graph (entities, typed relations, evidence quotes, confidence scores), and Layer 2 is an \emph{epistemic super-domain layer} that annotates every relation with a categorical status using a rule engine combining regex patterns over the verbatim evidence text with document-type inference. The two layers are produced in the same pipeline, persisted in the same artifact (\texttt{claims\_layer.json}), and rendered together in the visualization. A separate analyst-persona stage consumes the classified graph (not the source text) to produce a proactive briefing -- contradictions surfaced, prophetic claims separated from asserted ones, follow-up actions recommended.

Two design choices distinguish the framework from prior work. The first is \emph{domain pluggability}: the pipeline is domain-agnostic, and a domain is defined by four small configuration files (schema in YAML, extraction prompt in Markdown, epistemic rules in Python, analyst persona in YAML). The framework ships with four pre-built domains -- drug-discovery, contracts, clinicaltrials, and fda-product-labels -- and adding a new one requires no Python changes to the core pipeline. The second is the \emph{persona dual-use pattern}: a single YAML string per domain (the analyst persona) is the system prompt for both a reactive chat surface (workbench, where users ask ad-hoc questions) and a proactive briefing surface (the auto-generated narrator output). The two surfaces stay coherent because they derive from the same persona.

This paper is the second in a sequence. The first~\cite{bhatt2026epistract_v1} described the v1 framework -- a single-domain (drug-discovery) system validated against six scenarios -- as a refutation of similarity-based GraphRAG~\cite{bhatt2025latentkg, edge2024graphrag} for biomedical literature analysis. The current paper documents the framework's evolution into v3.2: a two-layer architecture, four pre-built domains, the persona dual-use pattern, and the seventh validation scenario (clinicaltrials). The v1 paper's argument -- that LLM \emph{comprehension} produces typed, evidence-backed relations that embedding-similarity cannot -- remains intact and is the foundation Layer 1 still rests on. What is new is Layer 2 and the framework around it.

The remainder of this paper is organized as follows. Section~\ref{sec:architecture} describes the two-layer architecture and the pipeline. Section~\ref{sec:epistemic} details the epistemic super-domain layer -- the rule engine, the status taxonomy, and the persona-driven briefing. Section~\ref{sec:domains} describes the four pre-built domains and the \texttt{/epistract:domain} wizard for adding new ones. Section~\ref{sec:evaluation} reports validation across all seven scenarios spanning two active domains. Section~\ref{sec:evolution} compresses the GraphRAG $\to$ Epistract v1 $\to$ Epistract v3 trajectory. Sections~\ref{sec:collaboration}--\ref{sec:conclusion} discuss collaboration, availability, and limitations.
